%    \begin{macrocode}
%<*exemple>

\newif\ifpdf
\ifx\pdfoutput\undefined
  \pdffalse
\else
  \pdfoutput=1
  \pdftrue
\fi

\ifpdf
\documentclass[10pt,a4paper,pdftex,french]{report}
\else
\documentclass[10pt,a4paper,french]{report}
\fi
\let\ifpdf\undefined

\title{Document test du package \texttt{varchap}}
\author{Sébastien \sc Le Callonnec}
\date{\today}

\usepackage{verbatim}
\usepackage{times}
\usepackage{amssymb}
\usepackage[msg]{varchap}
\usepackage[T1]{fontenc}
\usepackage{babel}
\usepackage{makeidx}\makeindex

\begin{document}
\maketitle

\renewcommand{\contentsname}{Sommaire}
\newcommand{\lestyle}[1]{\textsf{#1}\index{#1@\textsf{#1}}}
\newcommand{\cmd}[1]{\texttt{\symbol{92}#1}\index{#1@\texttt{\symbol{92}#1}}}
\newcommand{\cmdx}[1]{\index{#1@\texttt{\symbol{92}#1}}}
\newcommand{\corps}{%
Le but de ce document est de fournir une présentation du package 
\texttt{varchap} et de ses différentes possibilités.

Ce chapitre utilise le style <<~\lestyle{\nomstyle}~>>.}


\newcommand{\imprimechapitre}{%
\chapter{Présentation\\style \nomstyle}
\corps
}
\definecolor{rouge}{rgb}{0.6,0,0}
\definecolor{tresrouge}{rgb}{1,0,0}
\definecolor{violet}{rgb}{0.6,0,1}

\tableofcontents

\setlength{\parskip}{6pt}

\chapter{Introduction}
Ce document est le fichier de test du package \texttt{varchap}. C'est
de plus un manuel d'utilisation de ce package, fournissant des commandes
à utiliser afin de personnaliser l'apparence des styles.

Pour toute remarque ou question, écrire à~:

\begin{center}
Sébastien \textsc{Le Callonnec}\\
email:<\texttt{sebastien@weblogism.com}>
\end{center}

\imprimechapitre

Ce style est le style par défaut. On peut y accéder en appelant la
commande \cmd{getEllen}.

\coulNumChap{}
\imprimechapitre

Changement de couleur du numéro. Ce résultat a été obtenu en utilisant
la commande~:


\cmd{coulNumChap}

\chapLong{-40mm}{40mm}{}
\tailleTitreChap{\Huge\itshape}
\makeatletter
\fonteNumChap{\@fontenumer@{T1}{phv}}
\makeatother
\imprimechapitre

Ce style a été obtenu à partir du style de base \lestyle{Ellen}, en utilisant les
commandes~:


\begin{verbatim}
\chapLong{-40mm}{40mm}{}
\tailleTitreChap{\Huge\itshape}
\makeatletter
\fonteNumChap{\@fontenumer@{T1}{phv}}
\makeatother
\end{verbatim}


\getChristelle
\imprimechapitre

Le style \lestyle{Christelle} est accessible par l'option \texttt{christelle}
ou la commande \cmd{getChristelle}.


\getAudrey
\imprimechapitre

L'exemple du style \lestyle{Audrey} montre que l'on peut insérer des filets 
dans l'interstice entre le numéro de chapitre et le titre. Ainsi, la commande
à utiliser est~:

 \verb+\insereTampon{\vspace{-40mm}\rule{\linewidth}{0.4pt}}+ 
 
\noindent pour avoir un filet tel qu'il est sur cette page. Il est ensuite possible de 
jouer avec les écarts entre numéro et filet, entre filet et titre, et la 
longueur/épaisseur du filet. Voire même en mettre plusieurs.

Le style \lestyle{Audrey} est accessible par l'option \texttt{audrey}
ou la commande \verb+\getAudrey+.



\getChristelle
\coulTitreChap{\color{rouge}}
\imprimechapitre

Le style \lestyle{Christelle} a été utilisé ici avec un changement de couleur 
du numéro de chapitre. La commande employée est~:

\verb+\coulTitreChap{\color{rouge}}+

\noindent L'interface utilise ainsi le package \texttt{color} afin de changer
les couleurs. Il faut toutefois définir les couleurs au préalable avec 
\verb+\definecolor+. Par exemple, ici~:

\verb+\definecolor{rouge}{rgb}{0.6,0,0}+

Pour plus d'informations, cf. doc du package \texttt{color}.

Le package \texttt{varchap} définit les couleurs \texttt{grisaille} et 
\texttt{grisaillepale} qui correspondent aux différents tons de gris des
numéros de chapitre présents dans ce document.


\getUndine{70pt}{3}
\imprimechapitre

Le style \lestyle{Undine} utilise le package \texttt{graphicx} pour afficher
le mot <<~Chapitre~>> le long du numéro de chapitre, en rotation de 90°.

\textbf{Attention!} Cette rotation n'apparaît que dans le \textsc{ps}, et pas
dans le \textsc{dvi}. Donc, pas de panique si vous ne la voyez pas dans le 
\textsc{dvi}.

Le style \lestyle{Undine} est accessible par l'option \texttt{undine} ou par
la commande~: 

\cmd{getUndine}\verb+{+\textit{longueur}\verb+}{+\textit{nombre}\verb+}+

\noindent où \textit{longueur} est la longueur de l'indentation du titre du
chapitre, et \textit{nombre} le nombre de ligne de retour en arrière.

Par défaut, l'option \texttt{undine} exécute la commande \verb+\getUndine{0pt}{2}+,
ce qui donne le résultat de la page suivante.

Ici, la commande est~: 

\verb+\getUndine{70pt}{3}+

\getUndine{0pt}{2}
\imprimechapitre

\getUndine{0pt}{2}
\renewcommand{\chaptername}{Quelque chose}
\imprimechapitre
\renewcommand{\chaptername}{Chapitre}

Renommage de <<~Chapitre~>> en <<~Quelque chose~>> par~:

\verb+\renewcommand{\chaptername}{Quelque chose}+


\getUndine{0pt}{2}
\coulNumChap{}
\coulTitreChap{\color{violet}}
\imprimechapitre

Une nouvelle fois le style \lestyle{Undine}. Cette fois, on passe à un
numéro de chapitre en noir, et à un titre en violet. Le code~:

\begin{verbatim}
\getUndine{0pt}{2}
\definecolor{violet}{rgb}{0.6,0,1}
\coulNumChap{}
\coulTitreChap{\color{violet}}
\end{verbatim}

\getLisbeth
\imprimechapitre

Le style \lestyle{lisbeth} centre tout et fait se chevaucher le numéro
et le titre. Tout ça est modulable grâce à \cmd{intersticeskip}.


\getUndine{0pt}{2}
\chapter{Et maintenant un long titre de chapitre pour voir}
\corps


\getUndine{0pt}{2}
\chapter*{Version étoilée}
\corps

La version étoilée de \verb+\chapter+,\cmd{chapter*}, donne des choses pas trop belles.
Il faut alors manipuler les distances pour arranger les choses.


\getClaudia
\imprimechapitre

Le style \lestyle{Claudia} utilise lui aussi le <<~tampon~>> pour mettre
l'étoile et les filets. Vous pouvez redéfinir \verb+\monetoile+ pour obtenir
un autre symbole (par exemple 

\verb+\renewcommand{\monetoile}{$\spadesuit$}+

\noindent me vient
à l'esprit), ou tout simplement redéfinir \cmd{insereTampon} pour complètement
changer.


\getUndine{0pt}{2}
\preChapitre{\vspace*{-3cm}%
{\raggedleft{\itshape{%
L'\'Eternité, ce sera peut-être long,\\
mais on n'aura qu'à couper à la moitié.}}\\[1.2\baselineskip]
-- Un sage\par}%
}
\imprimechapitre

Cette fois, on remplit le \cmd{prechapitre} avec une citation. La mise 
en page est à étudier à partir du code suivant~:

\begin{verbatim}
\preChapitre{\vspace*{-3cm}%
{\raggedleft{\itshape{%
L'\'Eternité, ce sera peut-être long,\\
mais on n'aura qu'à couper à la moitié.}}\\[1.2\baselineskip]
-- Un sage\par}%
}
\end{verbatim}


\getSylvia
\imprimechapitre


\getClaudia
\insereTampon{}
\makeatletter
\renewcommand{\@chapitrenom}{\color{tresrouge}\MakeUppercase{\chaptername}}
\makeatother
\coulTitreChap{\color{tresrouge}}
\coulNumChap{\color{tresrouge}}
\chapLong{-2cm}{}{}
\renewcommand{\prechapitre}{%
\boiteChapitre{5cm}{7cm}{-4cm}{
	  \typeout{Ô, le bo dessin^^J%
	  Ich bin ein Bild.}%
	  \parbox{5cm}{\scriptsize\textbf{Ô, le bo dessin\\Ich bin ein Bild.}}%
	}}
\imprimechapitre

Pour finir, on remplit notre boîte du \cmd{prechapitre} par une image \textsc{eps} 
ou \textsc{jpg} selon que  l'on fasse un document \textsc{pdf} ou PostScript. Le code~:

\begin{verbatim}
\getClaudia
\insereTampon{}
\makeatletter
\renewcommand{\@chapitrenom}{\color{tresrouge}%
\MakeUppercase{\chaptername}}
\makeatother
\coulTitreChap{\color{tresrouge}}
\coulNumChap{\color{tresrouge}}
\chapLong{-2cm}{}{}
\preChapitre{%
\boiteChapitre{5cm}{7cm}{-4cm}{
	\includegraphics[width=12cm]{image.png}}
\end{verbatim}

\printindex
\end{document}
%</exemple>
%    \end{macrocode}
%%% Local Variables: 
%%% mode: latex
%%% TeX-master: t
%%% End: 
